\uuid{snqM}
\titre{Partie entière}
\auteur{bodin ; hanani ; mzari}
\organisation{exo7}
\licence{CC BY-NC-SA 4.0}
\datecreate{2020-01-20}
\niveau{L1}
\module{Analyse}
\chapitre{Propriétés de R}
\sousChapitre{Propriétés des nombres réels}
\difficulte{2}
\theme{Propriétés de nombres réels, partie entière}
\exo7qcmid{lille-1.reels.013}

\contenu{
\begin{question}
\qclassification{120.03}
\qtags{motcle=partie entière}
Pour $x\in \Rr$, on définit $G(x) = E(10x)$.
\begin{answers}
    \bad{$G(\frac23) = 66$}

    \bad{$\forall x>0 \quad G(x) \ge 1$}

    \bad{$G(x)=10 \iff x \in\{10,11,12,\ldots,19\}$}

    \good{$G(x)=G(y) \implies |x-y| \le \frac{1}{10}$}

\end{answers}
\begin{explanations}
La fonction $G$ est très similaire à la fonction partie entière.
\end{explanations}
\end{question}
}
